\chapter{Conductors and Capacitors} \label{chap:conductors-and-capacitors}

\section{Conductors}

\section{Capacitors} \label{sec:capacitors}

\textbf{Capcitors} are devices that store electric charge.
Unlike batteries, they are purely electrical rather than electrochemical.
Physically,
they are composed of two conductors separated by a distance,
which have opposite charges of equal magnitude when charged.

How well a capacitor stores charge is its \textbf{capacitance} (\(C\)),
with SI units of farads \unit{\farad}.
The capacitance is defined as
\begin{equation}
    C = \frac{q}{V},
\end{equation}
where \(q\) is the magnitude of charge on each conductor
and \(V\) is the voltage between the two conductors.

\subsection{Parallel Plate Capacitors}

\textbf{Parallel plate capacitors} are the most common type of capacitor.
They are composed of two parallel plates,
and their capacitance \(C\) is given by
\begin{equation}
    C = \frac{\epsilon_0 \kappa \mathrm{d} A}{d},
\end{equation}
where \(\kappa\) is the unitless constant of the \textbf{dielectric},
\(A\) is the area of each plate,
and \(d\) is the distance between the plates.

A dielectric is a insulator inserted between the plates of the capacitors
to increase the capacitance.
A vacuum will have a constant of \(\kappa = 1\),
since \(\epsilon_0\) is the vacuum permittivity after all.
Air has a dielectric constant of slightly less than 1.
When the voltage across the capacitor is too high,
\textbf{dielectric breakdown} occurs;
the voltage at which this occurs is dependant on the dielectric.
